\section{Coxeter Groups and the Kaleidoscope}

The previous section enumerated the family of regular hyperbolic honeycombs and showed why the list is finite. This section steps back to the idea underneath that list. There is one idea, and it is a kaleidoscope. A few mirrors, set at exact angles, reflect a single small wedge into a whole symmetric world.

\begin{principle}[The kaleidoscope]
A handful of mirrors generates a universe. The mesh is one chamber reflected without end.
\end{principle}

This is the deepest structure in the theory. It is not only the mesh. The same idea gives the $24$ sites, the symmetry that carries spin, and the gauge structure. One small picture, held clearly, makes all of the later machinery feel inevitable. The space, the directions out of a dock, and the law those directions obey all come out of the same few mirrors.

\subsection{The kaleidoscope}

A child's kaleidoscope is two or three strips of mirror in a tube. Drop in a few colored beads and look down the tube. You do not see a few beads. You see a perfect six-fold rosette, repeated outward, filling the circle.

Nothing was added. The mirrors took one small wedge of reality, reflected it, then reflected the reflections, until the wedge tiled the whole disk. That is the entire engine.

\begin{itemize}
\item Take one small region, a \emph{chamber}, bounded by mirror walls.
\item Reflect the chamber across each wall.
\item Reflect the copies across their walls. Forever.
\end{itemize}

The copies fit together with no gaps and no overlaps. They tile the space. The wedge you started with is the only original. Everything else is its mirror image, or the image of an image. Repeated reflection does not decorate the space. It \textbf{builds} it.

The mesh of the universe is exactly this. One chamber. A few mirrors. The chamber reflected endlessly into the growing honeycomb of docks.

\subsection{What a Coxeter group is}

A \emph{Coxeter group} is the group of all those reflections and all their combinations. Start with a few mirrors. Each mirror is a reflection, a flip across a wall. The group is everything you can reach by flipping, then flipping again, in any order, as many times as you like.

It is generated by the mirrors and almost nothing else. There are only two kinds of relation.

\begin{itemize}
\item \textbf{A mirror is its own inverse.} Reflect across one wall, then reflect across the same wall, and you are back where you started. Flip twice, no change. In symbols $r_i^2 = 1$.
\item \textbf{Two mirrors that meet.} If two mirrors meet at an angle of $\pi / m$, then alternating between the two reflections is a rotation that returns to the start after $m$ rounds. In symbols $(r_i r_j)^{m_{ij}} = 1$.
\end{itemize}

That is all. No other relations. Write down which mirrors there are, and for each pair the single integer $m$, and the entire infinite group is fixed. The tessellation is fixed with it.

\begin{principle}[One integer per pair]
One integer per pair of mirrors is the whole input. The chamber, the group, and the endless tiling all follow.
\end{principle}

The bookkeeping is tiny. A small diagram, one node per mirror and one labeled edge per meeting pair, records every integer. This is the \emph{Coxeter diagram}, and the matching \emph{Schl\"afli symbol} like $\{3,4,3,4\}$ is the same data laid out in a row. The next section reads this symbol as a recipe for unfolding one chamber into the tiling.

\begin{definition}[Coxeter group]
A Coxeter group is the group generated by a finite set of reflections $r_1, \dots, r_n$ subject only to the relations $r_i^2 = 1$ and $(r_i r_j)^{m_{ij}} = 1$, where each $m_{ij}$ is a whole number recording that mirrors $i$ and $j$ meet at angle $\pi / m_{ij}$. The integers $m_{ij}$ are the complete description of the group and of the tiling it generates.
\end{definition}

\subsection{Why the angles must be exact}

The mirror angles are not free. They must be exact integer submultiples of $\pi$. The angle between two mirrors must be $\pi$ divided by a whole number. So $\pi/2$, $\pi/3$, $\pi/4$, $\pi/5$, and so on. Never $\pi$ over a fraction.

The reason is the tiling itself. Put two mirrors at an angle and reflect the wedge between them around and around. The copies sweep through a full turn of $2\pi$. For the copies to meet up cleanly when they come all the way around, the wedge angle has to divide the full turn a whole number of times.

If the angle is $\pi/3$, then six copies of the wedge fill the circle and the last one lines up exactly with the first. Clean. If the angle is anything else, the copies drift. The last one lands slightly off the first. They overlap, or they leave a gap. The tiling fails.

\begin{principle}[Integer angles tile]
Integer angles are the condition for the reflected images to fit. Anything else tears or doubles the fabric.
\end{principle}

This single condition is why the honeycombs form a finite, enumerated family and not a free-for-all. Only certain integer combinations close up in a given curvature. The previous section walked that enumeration and showed exactly where the list runs out. The kaleidoscope explains \emph{why} there is a list at all. The mirrors are only allowed to meet at the angles that tile.

\subsection{The three regimes}

The same kaleidoscope idea runs in three settings, set by how the wedge angles add up. This is the classification of all Coxeter groups, and it is just the three curvatures of the geometry section.

\begin{table}[ht]
\centering
\begin{tabular}{@{}llll@{}}
\toprule
regime & the angles & the group & what it tiles \\
\midrule
spherical & sum to more than flat & finite & a sphere, closing into a polytope \\
Euclidean & sum to exactly flat & affine, infinite, flat & flat space, a periodic lattice \\
hyperbolic & sum to less than flat & infinite, negatively curved & hyperbolic space, ever opening \\
\bottomrule
\end{tabular}
\caption{The three Coxeter regimes, set by how the chamber angles compare to flat. They match the three constant-curvature geometries one to one.}
\label{tab:three-regimes}
\end{table}

\textbf{Spherical.} The mirrors meet so tightly that the reflections close into a finite group. The wedge wraps around a sphere a whole number of times and stops. The output is a regular polytope, like the cube or the $24$-cell. Finitely many chambers, then done.

\textbf{Euclidean.} The angles balance on the knife edge of flatness. The reflections include translations, and the wedge marches across flat space forever in a repeating grid. This is the affine case, the parabolic edge between the other two. It is where the flat cusp of the mesh lives.

\textbf{Hyperbolic.} The angles fall short of flat. There is always more room. The reflections never close and never settle into a periodic grid. The wedge tiles a space that opens faster than flat space ever could. This is the regime the substrate lives in.

The mesh is a hyperbolic kaleidoscope whose chamber, by one balanced choice of angles, carries a flat Euclidean cusp at its boundary. The bulk is the hyperbolic regime. The cusp is the affine regime sitting inside it. A later section owns the detail of how the affine cubic layer hides inside the hyperbolic group.

\subsection{One idea, three things at once}

Here is why the kaleidoscope is the spine of the whole theory and not only a fact about space. The same reflection-group machinery that builds the mesh also builds the directions and the symmetry. One choice of mirrors hands you all three.

\begin{itemize}
\item \textbf{The space.} Reflect the chamber and you get the honeycomb. This is the mesh, the growing crystal of docks.
\item \textbf{The directions.} The mirrors of certain groups also carve out a \emph{root system}, a fixed set of vectors closed under the reflections. For the mesh's chamber that root system is the $D_4$ roots, the $24$ vectors that form the $24$-cell. These are the $24$ \textbf{sites}, the $24$ directions out of a dock. The same mirrors that tile the space also fix where a dock can point.
\item \textbf{The symmetry.} The full group of a finite root system is its \emph{Weyl group}, the exact symmetry of those directions. For the mesh's cell the relevant group is the exceptional $F_4$ symmetry of order $1152$. This is the symmetry the later physics rides on. It is what lets the sites carry spin and the gauge structure.
\end{itemize}

So one small kaleidoscope generates the room, the directions you can move in, and the law those directions obey. The space and the alphabet and the grammar all come out of the same few mirrors.

\begin{principle}[One input, three outputs]
A few mirrors fix the space, the $24$ directions, and the symmetry that carries spin. Not three inputs. One.
\end{principle}

This is the unity the theory keeps drawing on. When the spin section finds that the sites carry a spinor, it is reading a property of the same root system the mirrors built. The spin is not bolted on. It was already in the kaleidoscope, latent in the reflection group, waiting to be read off the directions.

\subsection{One structure, three roles}

The same $24$ sites play three roles at once. Not three structures bolted together. One structure, read three ways.

\begin{table}[ht]
\centering
\begin{tabular}{@{}ll@{}}
\toprule
role & what the $24$ sites are \\
\midrule
navigation & the address-tree generator, the bulk's branching, the growth that names every dock \\
symmetry & the $D_4$ root system, which carries spin and the gauge structure \\
physics & shared with the flat cubic cusp, the 3D space where physics lives \\
\bottomrule
\end{tabular}
\caption{The $24$ sites carry navigation, symmetry, and physics at the same time, with no trade-off between the roles.}
\label{tab:three-roles}
\end{table}

\begin{result}[No trade-off]
One structure carries navigation and symmetry and physics at once. There is no trade-off between the roles.
\end{result}

A grid that wins one role often loses another. A flat lattice that gives clean 3D physics carries no spinor. A pure 2D addressing grid carries no spin and no 3D space. The $\{3,4,3,4\}$ mesh refuses the trade. The very same $24$ directions that branch the address tree are the $D_4$ roots that carry spin, and they live on a mesh whose cusp is the flat 3D space of physics. One choice of mirrors, three roles, no compromise. A later section develops the navigation role in full.

\subsection{Why this is deep, not decorative}

It would be easy to take the kaleidoscope as a pretty way to draw a tiling. It is much more than that. The reflection-group viewpoint says that geometry is a group first. The picture is the group made visible. You do not place the docks and then notice a symmetry. The symmetry is primary, and the docks are its orbit.

This flips the usual order. Most lattices are drawn, then studied for whatever symmetry they happen to have. A Coxeter honeycomb is the reverse. The symmetry is the definition, and the lattice is forced. There is no freedom left to place anything by hand. Choose the mirrors and everything else is determined, the space, the cells, the directions, the law.

\begin{proposition}[Symmetry first, lattice forced]
In a Coxeter honeycomb the reflection group is the primitive object and the tiling is its orbit. Once the mirror angles are chosen, the space, the cells, the directions, and the law are all fixed with no further freedom. Nothing is placed by hand.
\end{proposition}

That forced, rigid, determined character is exactly what the theory wants from a substrate. A substrate placed by hand could have been placed differently, so it would beg the question of why this one. A substrate forced by a few integers is as close to inevitable as a structure can be. A later section develops this into the claim that the mesh is discovered, not invented.

\subsection{What is borrowed and what is new}

None of the mathematics here is the theory's invention. It is the deep and beautiful work of others.

\begin{itemize}
\item \textbf{H. S. M. Coxeter} unified regular polytopes, reflection groups, and the higher-dimensional and hyperbolic cases into one recursive language, the diagrams and the classification.
\item \textbf{Jacques Tits} built the algebraic theory of these groups and their buildings, making the word-and-chamber picture exact.
\item \textbf{Ernest Vinberg} developed the theory of reflection groups in hyperbolic space, the very regime the substrate lives in.
\end{itemize}

The root systems, the Weyl groups, the exceptional $F_4$, the affine and hyperbolic Coxeter groups, all of it is established mathematics, decades old and rock solid.

\begin{principle}[The single act of the theory]
The contribution of the theory is not the kaleidoscope. It is the choice of one specific kaleidoscope as the substrate of reality.
\end{principle}

Every other tessellation in the family is a kaleidoscope too. Each is a real, buildable world. The theory's single act is to pick one, the $\{3,4,3,4\}$ mesh, on the grounds that its mirrors alone deliver curvature, a flat 3D cusp, and a spinor on its $24$ sites all at once. The selection sections are the argument that this choice is forced. This section is only the idea the choice ranges over.

The mathematics is given. The reading of it as the floor of the world is the theory.
